% Chapter 2. Add new references to ../references.bib; keep claims tied to a citation.
\chapter{Background and literature review}
\label{ch:literature}

This chapter covers the background that the method relies on. \Cref{sec:lit-gqd,sec:lit-shell}
describe GQDs and the closed-shell condition that every design must meet. \Cref{sec:lit-sensing}
reviews gas sensing with graphene materials. \Cref{sec:lit-methods,sec:lit-ml} summarise the
quantum-chemical and machine-learning methods used for scoring. \Cref{sec:lit-repr,sec:lit-design,sec:lit-llm}
review molecular representations, computational molecular design and the use of LLMs in chemistry,
and \cref{sec:lit-gap} states the gap that this thesis addresses.

\section{Graphene and graphene quantum dots}
\label{sec:lit-gqd}

Graphene is a single layer of carbon atoms arranged in a honeycomb lattice \cite{novoselov2004}.
Each carbon is \(sp^2\) hybridised and contributes one \(p_z\) electron to a delocalised \(\pi\)
system. Infinite graphene has no band gap: its valence and conduction bands meet at the corners of
the Brillouin zone \cite{castroneto2009}. When graphene is cut into a flake a few nanometres wide,
the electrons are confined and a gap opens \cite{ponomarenko2008}. Such flakes are called graphene
quantum dots \cite{bacon2014gqd}. With hydrogen on every edge carbon, a GQD is chemically a large
polycyclic aromatic hydrocarbon (PAH), and in this thesis the two terms describe the same kind of
molecule.

GQDs are made in two ways. Top-down methods cut or oxidise larger carbon materials and give
mixtures of sizes and edge structures. Bottom-up methods build the flake from small organic
precursors and can give one exact structure \cite{yan2010large,narita2015nanographene}. Because
bottom-up synthesis can reach defined shapes, a computational design that names an exact
structure is a meaningful target.

The gap of a GQD depends on more than its size. Graphene has two kinds of straight edge, armchair
and zigzag. Zigzag edges carry electronic states localised at the edge with energies close to the
middle of the gap \cite{nakada1996edge}. A flake with long zigzag edges therefore has a smaller
gap and more reactive edge carbons than an armchair-edged flake of the same size. Shape also
controls whether the flake has unpaired electrons, which is the subject of the next section.

\section{Closed-shell and open-shell GQDs}
\label{sec:lit-shell}

The carbon atoms of a honeycomb lattice fall into two sublattices, A and B, and every
carbon-carbon bond joins an A atom to a B atom. The carbon skeleton of a GQD is therefore a
bipartite graph. A Kekul\'e structure, a pattern of alternating single and double bonds in which
every carbon has exactly one double bond, corresponds to a perfect matching of this graph. A perfect
matching can only exist if the two sublattices have the same number of atoms, \(N_A = N_B\)
\cite{cyvin1988kekule}. Equal numbers are necessary but not sufficient; some balanced benzenoids
still have no Kekul\'e structure \cite{cyvin1988kekule}.

When \(N_A \neq N_B\), the molecule cannot be closed-shell. For the half-filled Hubbard model on a
bipartite lattice, Lieb's theorem gives a ground-state spin of \(S = |N_A - N_B|/2\)
\cite{lieb1989}, in line with Ovchinnikov's earlier rule for alternant hydrocarbons
\cite{ovchinnikov1978}. Triangular flakes are the standard example. Their sublattice imbalance
grows with size and they have a magnetic ground state
\cite{fernandezrossier2007,wang2009topological}. Even among closed-shell structures, stability
differs; Clar's aromatic sextet theory relates it to the number of disjoint benzene-like sextets
that can be drawn \cite{clar1972sextet}.

For design, the consequence is simple. A GQD without a Kekul\'e structure is a radical. It is
reactive, hard to make and hard to describe with the closed-shell methods used for screening. The
workflow in this thesis therefore accepts only GQDs that have a Kekul\'e structure, and it checks
this condition with an exact matching algorithm (\cref{sec:hex}).

\section{Gas sensing with graphene materials}
\label{sec:lit-sensing}

In a chemiresistive sensor, an adsorbed gas changes the carrier density of the sensing layer and
therefore its resistance. \ce{NO2} is an electron acceptor. On graphene it withdraws electrons and
dopes the layer p-type, which is the effect that allowed Schedin et al.\ to detect single
molecules \cite{schedin2007detection}. First-principles calculations by Leenaerts et al.\ found that
\ce{H2O}, \ce{NH3}, CO, \ce{NO2} and NO all physisorb weakly on pristine graphene, that \ce{NO2}
is a strong acceptor, and that the amount of charge transferred depends on the orientation of the
molecule \cite{leenaerts2008adsorption}. Dopants and defects strengthen the binding of several
gases \cite{zhang2009dopants}, and reviews of graphene sensors describe the same balance between
sensitivity, selectivity and recovery \cite{yuan2013graphene}.

Recovery is the reverse problem. If a gas binds too strongly, it does not leave at room
temperature and the sensor stays saturated. Computational studies often estimate the recovery time
with transition-state theory as \(\tau = \nu^{-1}\exp(-\Eads/k_BT)\), with an attempt frequency
\(\nu\) of about \(10^{12}\)~s\(^{-1}\) \cite{peng2004cnt}. Because \(\tau\) depends exponentially on
the adsorption energy, a change of 0.1~eV in binding changes the recovery time by a factor of about
50 at room temperature. Peng et al.\ used such estimates to explain the slow recovery observed for
carbon-nanotube \ce{NO2} sensors \cite{peng2004cnt}.

GQDs have been combined with metal oxides in \ce{NO2} sensors. Lee et al.\ decorated \ce{SnO2}
nanodomes with GQDs and reported room-temperature detection with a limit of 1.1~ppb and good
selectivity against other gases \cite{lee2023gqdsno2}. In such studies the GQDs come from
solution synthesis and their exact structure is not controlled. We did not find studies that
search over exact GQD structures for \ce{NO2} sensing. \authornote{Confirm this with a systematic
search (Web of Science, Scopus, arXiv) before submission.}

\section{Computational methods for adsorption}
\label{sec:lit-methods}

\paragraph{Density functional theory.}
DFT is the standard method for adsorption energies and charge transfer in molecules of this size.
Its main weakness for this problem is delocalisation error: approximate functionals spread charge
too much, which favours fractional charge transfer and overbinds some radical and charge-transfer
complexes \cite{cohen2008insights}. Range-separated hybrid functionals reduce this error. Weak
binding of a gas to an aromatic surface is dominated by dispersion, so the functional also needs a
dispersion term. \(\omega\)B97X-V is a range-separated hybrid that includes the VV10 nonlocal
correlation functional for dispersion \cite{mardirossian2014wb97xv,vydrov2010vv10}, and it performs
well in broad benchmarks of molecular energies \cite{bursch2022bestpractice}. Small basis sets such
as def2-SVP \cite{weigend2005def2} make large systems affordable but cause basis set superposition
error, which overbinds weak complexes unless a counterpoise correction is applied
\cite{boys1970counterpoise}. Atomic charges are not observables, and their values depend on the
partitioning scheme. In this work DFT charges are computed with meta-L\"owdin populations
\cite{sun2014metalowdin}.

\paragraph{Extended tight binding.}
GFN2-xTB is a semiempirical tight-binding method parametrised for most of the periodic table. It
includes multipole electrostatics and a density-dependent dispersion correction
\cite{bannwarth2019gfn2,bannwarth2021xtb}. It is orders of magnitude faster than hybrid DFT and
gives good geometries for organic molecules, which makes it suitable for screening. Its
self-consistent charges come from a DFT-like model, so it is expected to share the tendency of
simple functionals to over-delocalise charge. \Cref{sec:res-dft-small} tests this directly for
\ce{NO2} and \ce{O2} on GQDs.

\paragraph{GPU acceleration.}
PySCF is an open-source Python framework for quantum chemistry
\cite{sun2018pyscf,sun2020pyscf}. Its GPU extension, GPU4PySCF, computes the electron-repulsion
integrals on the GPU; its authors report a speedup of about two orders of magnitude over
multithreaded CPU Hartree-Fock on one NVIDIA A100 \cite{li2025gpu4pyscf}. This speedup is what makes
the DFT benchmark in \cref{sec:e4} possible within a day.

\section{Machine-learned property prediction}
\label{sec:lit-ml}

Fast predictors of quantum-chemical properties are often used to filter candidates before an
expensive calculation. For conjugated hydrocarbons, H\"uckel theory \cite{huckel1931} gives a
cheap estimate of the frontier orbital energies from the connectivity alone, and it is a natural
physical baseline. Gradient-boosted trees such as XGBoost and LightGBM \cite{chen2016xgboost,ke2017lightgbm}
work well on hand-made descriptors, and their hyperparameters can be tuned automatically, for
example with Optuna \cite{akiba2019optuna}. Graph neural networks learn directly from the molecular
graph \cite{gilmer2017mpnn}; the graph isomorphism network (GIN) is a simple and expressive variant
\cite{xu2019gin}. Large datasets of PAHs exist. The COMPAS project, for example, published about
34{,}000 cata-condensed PAHs computed with GFN2-xTB and about 9{,}000 with DFT, and compared the two
levels of theory \cite{wahab2022compas}. Our predictor (\cref{sec:predictor}) uses a H\"uckel
baseline, physics-based features, boosted trees and a GIN-style network.

\section{Molecular representations for generation}
\label{sec:lit-repr}

Most generative methods write molecules as text or as graphs. SMILES encodes a molecule as a
string of atoms, bonds, branches and ring-closure digits \cite{weininger1988smiles}. A single wrong
character can make a SMILES string invalid, and large fused aromatic systems need many ring-closure
digits that must be paired correctly. SELFIES was designed so that every string decodes to a
molecule with valid valences \cite{krenn2020selfies}. Valid valences are not enough for GQDs,
however: a SELFIES string can still decode to an open-shell or non-planar structure.

For benzenoid hydrocarbons there is a more natural description. The carbon skeleton is fully
defined by the set of hexagons it occupies on the honeycomb lattice, a structure known in chemical
graph theory as a polyhex \cite{cyvin1988kekule}. In this description, connectivity, size and
shape can be read directly from the list of hexagons, and a local change to the shape is a local
change to the list. \Cref{sec:hex} uses this description as the language between the LLM and the
code.

\section{Computational molecular design}
\label{sec:lit-design}

Inverse design starts from a target property and searches for structures that have it
\cite{sanchezlengeling2018inverse}. Three families of methods are common. Generative models such as
variational autoencoders map molecules to a continuous space in which optimisation is easier
\cite{gomezbombarelli2018vae}. Reinforcement learning trains a generator to maximise a reward;
REINVENT is a widely used example \cite{olivecrona2017reinvent}. Genetic algorithms keep a
population of candidates and improve it by mutation, crossover and selection; graph-based genetic
algorithms are simple and strong baselines \cite{jensen2019graphga,nigam2020gad}.

Benchmarks compare these methods. GuacaMol defined a set of goal-directed tasks
\cite{brown2019guacamol}. The PMO benchmark later limited every method to a fixed number of
calls to the objective function and showed that this budget changes the ranking of methods, with
genetic algorithms and REINVENT among the strongest \cite{gao2022pmo}. A fixed budget matters even
more in our setting, because each call to the objective is an xTB calculation of about a minute
rather than a cheap cheminformatics function.

Polycyclic aromatic systems and GQDs have their own design studies. GaUDI combined a diffusion
model with a property predictor to generate polycyclic aromatic systems with targeted electronic
properties \cite{weiss2023gaudi}. For GQDs, high-throughput time-dependent DFT on 284 structures
of different shape, size and doping, combined with machine learning, gave design rules for their
emission \cite{ozdemir2024gqd}, and Bayesian optimisation found the doped GQD with the strongest
absorption among 63 candidates with only 12 TDDFT calculations \cite{ozonder2022bo}. These studies
target optical properties, and none of them uses an LLM or addresses gas sensing.

\section{Large language models in chemistry}
\label{sec:lit-llm}

\paragraph{What LLMs can do.}
General-purpose LLMs answer many chemistry questions, but a benchmark of eight chemistry tasks
found that they struggle with tasks that need an exact understanding of molecular strings, such as
converting names to SMILES \cite{guo2023llmchem}. Fine-tuned LLMs can serve as property predictors
in low-data settings \cite{jablonka2024leveraging}. A recent review covers both uses and the
growing number of LLM agents in chemistry \cite{ramos2025review}.

\paragraph{Agents with tools.}
An LLM agent alternates between reasoning and calling external tools, and the tool results feed
back into the next step \cite{yao2023react}. ChemCrow connected an LLM to a set of expert chemistry
tools \cite{bran2024chemcrow}, and Coscientist planned and ran laboratory experiments with an LLM
in control \cite{boiko2023coscientist}. In both systems the chemistry is done by the tools, and
the LLM decides what to do next. El Agente applies the same pattern to computational chemistry: it
plans and runs quantum-chemistry workflows, including xTB and DFT calculations, from requests in
natural language \cite{zou2025elagente}. The agent in this thesis follows the same pattern
(\cref{sec:agent}).

\paragraph{LLMs as optimisers and mutation operators.}
LLMs can also act inside an optimisation loop. In OPRO the LLM sees earlier solutions with their
scores and proposes new ones \cite{yang2024opro}. Evolution through Large Models used an LLM as the
mutation operator for computer programs \cite{lehman2022elm}, and language-model crossover used
few-shot prompts to recombine parents \cite{meyerson2023lmx}. FunSearch paired an LLM with an
automatic evaluator and found new mathematical constructions \cite{romeraparedes2024funsearch}.
For molecules, MolLEO placed LLMs as mutation and crossover operators inside a graph-based genetic
algorithm and reported faster progress than the baseline algorithm on several drug-design tasks,
with molecules written as SMILES or SELFIES \cite{wang2024molleo}. Ojeda et al.\ evolved synthetic
routes instead of molecules: LLM agents act as mutation operators, and deterministic chemistry
tools carry out the routes, so the LLM cannot invent products \cite{ojeda2026harness}. This is the
same division of work as in this thesis. The idea has also reached inorganic materials. MatLLMSearch
used a pre-trained LLM inside an evolutionary loop to propose crystal structures
\cite{gan2025matllmsearch}, and LLEMA combined LLM proposals with chemistry-based rules and a memory
of past successes and failures for multi-objective materials discovery \cite{abhyankar2026llema}.

\paragraph{Fine-tuning.}
Logs of an optimisation loop can become training data. Low-rank adaptation (LoRA) fine-tunes a
large model by training small added matrices \cite{hu2022lora}, and direct preference optimisation
(DPO) trains a model to prefer one answer over another from pairs of examples
\cite{rafailov2023dpo}. Both are possible next steps for this work (\cref{sec:future}).

\section{Research gap}
\label{sec:lit-gap}

The work reviewed above leaves four gaps that this thesis addresses.
\begin{enumerate}
  \item LLM-guided evolution has been tested on drug-like molecules
        \cite{wang2024molleo,ojeda2026harness} and on crystalline materials
        \cite{gan2025matllmsearch,abhyankar2026llema}. We did not find a study that applies it to
        GQDs, nanographenes or other PAHs, where validity depends on a global condition (a
        Kekul\'e structure must exist).
  \item We did not find a study in which the LLM edits a geometric representation, here a hexagon
        grid, while code enforces the chemistry and explains every rejection.
  \item The objective here is a quantum-chemical score for one device application, \ce{NO2}
        sensing, and the screening method (xTB) is checked against DFT within the same work.
  \item The comparison with random mutation uses the same xTB budget, the same validation and
        selection, several seeds, and a complete log of every proposal.
\end{enumerate}
\authornote{A web search on 12 November 2026 found the studies cited in this section and no
direct precedent for the four points above. Confirm this with a database search (Web of Science,
Scopus, Google Scholar), and repeat it shortly before submission, because new work in this area
appears every month.}
